%! Mode:: "TeX:UTF-8"
%! TEX program = xelatex
\PassOptionsToPackage{quiet}{xeCJK}
\documentclass[withoutpreface,bwprint]{cumcm_thesis}

% ====== 基础工具 ======
\usepackage{etoolbox}
% \usepackage{xltabular}   % 已包含 longtable 和 tabularx
\usepackage{tabularx}
\BeforeBeginEnvironment{tabular}{\zihao{-5}}

% ====== 参考文献 ======
\usepackage[numbers,sort&compress]{natbib}

% ====== 数学与单位 ======
\usepackage{amsmath}
%\usepackage{siunitx}
\usepackage{xcolor}

% ====== 化学 ======
% \usepackage[version=4]{mhchem}

% ====== 图片与网页 ======
\usepackage{url}
\usepackage{subcaption}
\usepackage{graphicx}
\usepackage{textcomp}

% ====== 字体与中文 ======
\usepackage{fontspec}
\usepackage{xeCJK}
\setmainfont{TeX Gyre Termes}
% \setCJKmainfont{SimSun}

% ====== 表格 ======
\usepackage{array}
\usepackage{booktabs}
\usepackage{caption}
% \usepackage{makecell}

\newcolumntype{C}{>{\centering\arraybackslash}X}
\newcolumntype{R}{>{\raggedleft\arraybackslash}X}
\newcolumntype{L}{>{\raggedright\arraybackslash}X}

% ====== 论文信息 ======
\title{全国大学生数学建模竞赛论文模板}
\tihao{}
\baominghao{}
\schoolname{}
\membera{}
\memberb{}
\memberc{}
\supervisor{}
\yearinput{}
\monthinput{}
\dayinput{}



%%%%%%%%%%%%%%%%%%%%%%%%%%%%%%%%%%%%%%%%%%%%%%%%%%%%%%%%%%%%%
%% 正文
\begin{document}

\maketitle
\begin{abstract}
摘要

\textbf{对于问题一,}

\textbf{对于问题二,}

\textbf{对于问题三,}

\textbf{对于问题四,}

最后，



\keywords{关键词\quad  关键词\quad  关键词\quad  关键词 \quad 关键词}
\end{abstract}
%%%%%%%%%%%%%%%%%%%%%%%%%%%%%%%%%%%%%%%%%%%%%%%%%%%%%%%%%%%%% 

% \tableofcontents  % 目录
% \newpage

%%%%%%%%%%%%%%%%%%%%%%%%%%%%%%%%%%%%%%%%%%%%%%%%%%%%%%%%%%%%%  
\section{问题重述} % section:节，栏
\subsection{问题背景} % subsection：分段

中药材干燥是影响成品品质的重要工序，热风烘干是常用方式之一。该过程通常分为预热平衡与恒温干燥两个阶段，通过调节烘房温湿环境实现药材脱水。烘干时，工艺参数若选择不合理，会造成干燥效率低、能耗高、品质波动等问题；传统试验优化又存在成本高、周期长等局限。为了解决这些问题，本文将采用数理分析与数值仿真方法，研究中药材烘干过程中的变化规律。


%%%%%%%%%%%%%%%%%%%%%%%%%%%%%%%%%%%%%%%%%%%%%%%%%%%%%%%%%%%%% 

\subsection{问题要求}



\textbf{问题1}  
在热风烘干预热平衡阶段，考虑一形状近似为圆柱体、长度和半径已知的中药材。我们需要利用给出的药材初始温度、初始水分浓度、烘房温度与水分浓度随时间的变化情况，以及相关物性参数和水分浓度扩散系数经验公式，建立预热平衡阶段药材内部温度和水分浓度随时间、空间变化的数学模型，并基于该模型计算 给出的7个时刻下，距药材中心 \(0,0.5,1,1.5,2\mathrm{cm}\) 共5个位置处的药材温度和水分浓度，在论文中给出相应表格；同时计算 \(1800\mathrm{s}\) 内每隔 \(1\mathrm{s}\)、距药材中心每隔 \(0.1\mathrm{cm}\) 的完整结果，所有结果保留四位小数。

\textbf{问题2 }  
我们需要利用给出的相关经验公式，在预热平衡阶段与恒温干燥阶段的参数有所不同的前提下，建立整个烘干过程中药材温度和水分浓度变化规律的数学模型，并基于该模型计算 \(3\mathrm{h}\) 内每隔 \(0.5\mathrm{h}\)、距药材中心 \(0,0.5,1,1.5,2\mathrm{cm}\) 处的药材温度和水分浓度，按表3和表4的格式在论文中给出；同时计算每隔 \(1\mathrm{s}\)、距药材中心每隔 \(0.1\mathrm{cm}\) 的完整结果，所有结果保留四位小数。

\textbf{问题3}  
我们需要利用给出的水分浓度限值，确定药材烘干所需的时间（单位：\(\mathrm{h}\)），并计算每隔 \(6\mathrm{h}\)、距药材中心每隔 \(0.5\mathrm{cm}\) 的水分浓度，按表5的格式在论文中给出；同时计算药材内部水分浓度每隔 \(60\mathrm{s}\)、距药材中心每隔 \(0.1\mathrm{cm}\) 的完整结果，所有结果保留四位小数。

\textbf{问题4}  
在实际烘干过程中，药材会因水分流失发生尺寸变化。我们需要利用给出的药材尺寸变化数据和相关经验公式，确定药材的烘干时长，并计算每隔 \(6\mathrm{h}\)、距药材中心每隔 \(0.5\mathrm{cm}\) 的水分浓度，按表6的格式在论文中给出；同时计算药材内部水分浓度每隔 \(60\mathrm{s}\)、距药材中心每隔 \(0.1\mathrm{cm}\) 的完整结果，所有结果保留四位小数。





%%%%%%%%%%%%%%%%%%%%%%%%%%%%%%%%%%%%%%%%%%%%%%%%%%%%%%%%%%%%% 

\section{问题分析}
\subsection{问题一分析}

对于问题一，要求在预热平衡阶段，建立圆柱形中药材内部温度与水分浓度时空演化的数学模型，本质上是一个给定外部热湿环境边界后，求解药材内部非稳态传热传质耦合偏微分方程的初边值问题。

从物理机理看，中药材内部结构复杂，鉴于题目提供的密度、比热容、导热系数与有效扩散系数均为宏观等效系数，本文采取宏观连续介质假设，将药材划分为微元控制体，建立宏观连续介质模型。同时，针对圆柱形及初始分布均匀的特点，忽略轴向与圆周方向的变化，简化为\textbf{一维径向模型}。基于傅里叶导热定律与有效Fick定律建立控制方程。由于水分扩散系数 \(D\) 是水分浓度 \(C\) 与温度 \(T\) 的非线性函数，模型呈现强非线性特征。边界条件方面，药材中心满足对称条件，表面则发生对流换热与对流传质，属于第三类边界条件。

在模型求解上，由于难以获得解析解，本文采用\textbf{显式有限差分法（FDM）}进行数值求解。求解主要面临两大难点：一是环境数据离散，且显式差分对边界条件的一阶导数极为敏感，常规线性插值易引发数值震荡，故引入\textbf{PCHIP插值}，保证插值曲线一阶导数连续且单调，避免非物理的过冲现象；二是需妥善处理圆柱中心的坐标奇点及表面通量平衡，为此本文利用洛必达法则处理中心节点，并基于有限体积法思想，保障数值计算的稳定性与物理守恒性。最后，从全局能量守恒与水分质量守恒两个维度对算法进行验证，确保结果可靠并满足精度要求。




\subsection{问题二分析}	
对于问题二，

\subsection{问题三分析}
对于问题三，

\subsection{问题四分析}
对于问题四，

%%%%%%%%%%%%%%%%%%%%%%%%%%%%%%%%%%%%%%%%%%%%%%%%%%%%%%%%%%%%% 

\section{模型假设}

为简化问题，本文做出以下假设：

\begin{itemize}[itemindent=2em] %创建一个由点开头的无序列表
\item 假设1：药材圆柱模型下上下底面受热对中心无影响
\item 假设2
\item 假设3
\end{itemize}

%%%%%%%%%%%%%%%%%%%%%%%%%%%%%%%%%%%%%%%%%%%%%%%%%%%%%%%%%%%%% 

\begin{table}[H]
\centering
\begin{tabularx}{\textwidth}{CLC}
\toprule
符号 & 说明 & 单位 \\
\midrule
\(t\)      & 时间                                   & \(\mathrm{s}\)  \\
\(L\)      & 药材长度                               & \(\mathrm{m}\) \\
\(r\)      & 到药材中心的距离                       & \(\mathrm{m}\) \\
\(R\)      & 药材半径                               & \(\mathrm{m}\) \\
\(T_0\)    & 药材初始温度                           & \(^{\circ}\mathrm{C}\) \\
\(T(r,t)\) & 药材内部温度                           & \(^{\circ}\mathrm{C}\) 或 \(\mathrm{K}\) \\
\(T_s\)    & 药材表面温度                           & \(^{\circ}\mathrm{C}\) \\
\(T_\infty(t)\) & 烘房温度                     & \(^{\circ}\mathrm{C}\) \\
\(C_0\)    & 药材初始水分浓度                       & \(\mathrm{kg/kg}\) \\
\(C(r,t)\) & 药材水分浓度             & \(\mathrm{kg/kg}\) \\
\(C_s\)    & 药材初始水分浓度                       & \(\mathrm{kg/kg}\) \\
\(C_\infty(t)\) & 烘房空气水分浓度             & \(\mathrm{kg/kg}\) \\
\(\rho\)   & 药材密度                               & \(\mathrm{kg/m^3}\) \\
\(c_p\)    & 药材比热容                             & \(\mathrm{J/(kg\cdot K)}\) \\
\(k\)      & 药材热传导系数                         & \(\mathrm{W/(m\cdot K)}\) \\
\(\alpha\) & 热扩散系数（\(\alpha = k/(\rho c_p)\)） & \(\mathrm{m^2/s}\) \\
\(h\)      & 对流换热系数                           & \(\mathrm{W/(m^2\cdot K)}\) \\
\(h_m\)    & 对流传质系数                           & \(\mathrm{m/s}\) \\
\(D\)      & 水分浓度扩散系数                       & \(\mathrm{m^2/s}\) \\
\(Q_v\)    & 单位体积热量累积率                     & \(\mathrm{W/m^3}\) \\
\(q\)      & 热流密度矢量                           & \(\mathrm{W/m^2}\) \\
\(J_c\)    & 水分扩散通量矢量                       & \(\mathrm{kg/(m^2\cdot s)}\) \\
\(N_r\)    & 空间径向节点总数                       & --- \\
\(N_t\)    & 时间节点总数                           & --- \\

\bottomrule
\end{tabularx}
\label{tab:符号说明}
\end{table}



%%%%%%%%%%%%%%%%%%%%%%%%%%%%%%%%%%%%%%%%%%%%%%%%%%%%%%%%%%%%% 

\section{问题一的模型的建立和求解}
\subsection{模型选取}

在建立中药材热风干燥数学模型前，需明确描述传热传质过程的数值模拟方法。根据现有研究，农产品热风干燥的数值模拟方法主要可分为三类：干燥动力学模型、连续介质假设模型和孔道网络模型。

干燥动力学模型以薄层干燥为对象，通过水分比等指标描述平均含水率变化。其优点是形式简单、计算量小，能快速预测干燥特性；但多依赖经验参数，物理机理不明确，难以刻画内部温度与水分浓度的空间分布。

连续介质假设模型将中药材视为等效多孔介质，基于质量、能量和动量守恒建立偏微分控制方程。其优点是能从宏观尺度描述温度场与水分浓度场的时空演化，物理意义清晰，计算效率较高，是农产品干燥研究中最常用的方法。但该方法基于连续性假设，对具有各向异性或复杂孔隙结构的物料描述不够精细。

孔道网络模型从物料内部孔隙结构出发，将多孔介质离散为节点与孔道网络，通过毛细压力、蒸发冷凝等机制模拟水分迁移。其优点是能反映真实孔隙结构的影响，可解释“湿斑”和“不规则前沿”等现象；但建模复杂、计算量大，三维模拟尚不成熟，难以直接应用于工程尺度预测。

综上，三类方法各有侧重。考虑到题目给出的密度、比热容、导热系数与有效扩散系数均为宏观等效系数，且问题需要获取药材内部温度与水分浓度的时空变化规律及烘干时长，本文选择连续介质假设模型进行模拟。即将药材等效为均匀连续的多孔介质，建立描述其内部水分浓度与温度时空演化的偏微分方程组，结合给定的初始条件与边界条件，使用数值方法求解，最终得到温度和水分浓度分布。
\begin{figure}
    \centering
    \includegraphics[width=1\linewidth]{figures/selection.pdf}
    \caption{三种热风干燥数值模拟方法优劣势对比与选择依据}
    \label{fig:placeholder}
\end{figure}


\subsection{模型建立：热质传递控制方程的建立}

\subsubsection{热量守恒方程}

取药材内部一微元控制体，由能量守恒定律可知：单位时间内控制体内热量的累积等于净流入控制体的热量与内热源生成热量之和。本文忽略内热源，则单位体积内热量的累积率可表示为
\begin{equation}
Q_v = \rho c_p \frac{\partial T}{\partial t}
\end{equation}
式中，\(Q_v\) 为单位体积热量累积率，\(\rho\) 为药材密度，\(c_p\) 为药材比热容，\(T\) 为药材内部温度，\(t\) 为时间。

药材内部热量主要以导热方式传递。根据傅里叶定律，热流密度与温度梯度成正比、方向相反，即
\begin{equation}
q = -k \nabla T
\end{equation}
式中，\(q\) 为热流密度矢量，\(k\) 为药材热传导系数，\(\nabla\) 表示梯度算子，\(\nabla \cdot\) 表示散度算子。

由能量守恒定律，控制体内热量的累积率等于热流密度的散度，即
\begin{equation}
\rho c_p \frac{\partial T}{\partial t} = -\nabla \cdot q
\end{equation}
将式(2)代入式(3)，得到药材内部温度场的控制方程
\begin{equation}
\rho c_p \frac{\partial T}{\partial t} = \nabla \cdot (k \nabla T)
\end{equation}

\subsubsection{水分扩散方程}

药材内部水分的迁移以扩散为主。根据有效 Fick 定律，水分扩散通量与水分浓度梯度成正比、方向相反，即
\begin{equation}
J_c = -D \nabla C
\end{equation}
式中，\(J_c\) 为水分扩散通量矢量，\(D\) 为水分浓度扩散系数，\(C\) 为药材水分浓度。

由水分质量守恒定律，控制体内水分浓度的变化率等于水分扩散通量的散度，即
\begin{equation}
\frac{\partial C}{\partial t} = -\nabla \cdot J_c
\end{equation}
将式(5)代入式(6)，得到药材内部水分浓度场的控制方程
\begin{equation}
\frac{\partial C}{\partial t} = \nabla \cdot (D \nabla C)
\end{equation}

需要指出的是，题目附录中给出的扩散系数 \(D\) 并非常数，而是与水分浓度 \(C\) 及温度 \(T\) 相关的函数。因此，式(7)应进一步写为
\begin{equation}
\frac{\partial C}{\partial t} = \nabla \cdot \bigl( D(C,T) \, \nabla C \bigr)
\end{equation}

\subsubsection{一维径向模型}

考虑到药材初始温度与含水率分布均匀，且药材几何形状为圆柱体，可作如下简化：忽略沿轴向的变化，忽略沿圆周方向的变化，仅考虑从药材中心到表面的径向变化。由此，三维问题可退化为沿半径方向的一维问题。

在柱坐标系下，对任一标量场 \(u\) 和传输系数 \(q\)，有恒等式
\begin{equation}
\nabla \cdot (q \nabla u) = \frac{1}{r} \frac{\partial}{\partial r} \left( r q \frac{\partial u}{\partial r} \right)
\end{equation}
式中，\(r\) 为到药材中心的径向距离。

将式(9)分别代入温度控制方程(4)和水分浓度控制方程(8)，可得一维径向模型下的控制方程组：
\begin{equation}
\rho c_p \frac{\partial T}{\partial t} = \frac{1}{r} \frac{\partial}{\partial r} \left( r k \frac{\partial T}{\partial r} \right)
\end{equation}
\begin{equation}
\frac{\partial C}{\partial t} = \frac{1}{r} \frac{\partial}{\partial r} \left( r D(C,T) \frac{\partial C}{\partial r} \right)
\end{equation}

\subsection{初始条件与边界条件}

\subsubsection{初始条件}

烘干开始时，药材内部温度和水分浓度分布均匀，由题意可得
\begin{equation}
T(r,0) = 28^\circ\text{C}, \qquad C(r,0) = 2.55\,\text{kg/kg}, \qquad 0 \le r \le R
\end{equation}
式中，\(R\) 为药材半径。根据题目给定尺寸，半径 \(R = 0.02\,\text{m}\)。

\subsubsection{中心对称条件}

在圆柱中心 \(r=0\) 处，由对称性可知热流密度与水分通量均为零，即
\begin{equation}
\left. \frac{\partial T}{\partial r} \right|_{r=0} = 0, \qquad \left. \frac{\partial C}{\partial r} \right|_{r=0} = 0
\end{equation}

\subsubsection{表面对流换热条件}

烘房空气通过对流将热量传递给药材，药材表面温度不一定等于空气温度，药材内部导热与外部对流共同决定其升温过程。由于题目未直接指定药材表面温度，采用 Newton 对流边界条件：
\begin{equation}
-k \left. \frac{\partial T}{\partial r} \right|_{r=R} = h \left[ T(R,t) - T_f(t) \right]
\end{equation}
式中，\(h\) 为对流换热系数，\(T_f(t)\) 为烘房温度，其随时间的变化由附件1给出。

\subsubsection{表面对流传质条件}

药材表面含水率与空气水分浓度之间存在差异，水分先从药材内部扩散至表面，再通过气膜传递至空气中。因此，药材表面满足对流传质边界条件
\begin{equation}
-D \left. \frac{\partial C}{\partial r} \right|_{r=R} = h_m \left[ C(R,t) - C_\infty(t) \right]
\end{equation}
式中，\(h_m\) 为对流传质系数，\(C_\infty(t)\) 为烘房空气水分浓度，其随时间的变化由附件1给出。

\subsection{模型求解}

本文建立的一维径向传热传质控制方程属于非线性偏微分方程组，难以获得解析解。为此，本文采用显式有限差分法（FDM）对空间域和时间域进行离散，通过时间推进循环求解药材内部温度和水分浓度的时空分布。

\subsubsection{时空离散与网格划分}

在空间域上，将药材半径 \(R\) 离散为 \(N_r\) 个节点，空间步长 \(\Delta r = R/(N_r-1)\)，节点位置为 \(r_i = i\Delta r\)（\(i = 0,1,\dots,N_r-1\)）。在时间域上，将总时长离散为 \(N_t\) 个时刻，时间步长 \(\Delta t\)，记 \(t_n = n\Delta t\)（\(n = 0,1,\dots,N_t-1\)）。采用显式差分格式对控制方程进行离散，以保证算法的简洁性与计算效率。

\subsubsection{边界条件的数据预处理}

在算法的时间推进过程中，表面节点及内部节点的更新均高度依赖于烘房环境的实时边界值 \(T_\infty(t)\) 和 \(C_\infty(t)\)。由于仅提供了有限个离散时间点的环境参数数据，必须通过插值获取任意时间步 \(n+1\) 对应的环境值 \(T_\infty^{n+1}\) 和 \(C_\infty^{n+1}\)。

由于显式有限差分法对边界条件的一阶导数十分敏感，若采用常规线性插值，其在节点处导数不连续（折点）的非光滑特性极易引起数值震荡，进而导致求解发散。因此，本文选取 PCHIP（分段三次 Hermite 插值多项式）对烘房环境参数进行插值。PCHIP 不仅保证了插值曲线的一阶导数连续，符合物理量平滑变化的规律，同时\textbf{保持数据的单调性，避免非物理的过冲现象}。这使得插值得到的水分浓度不会出现负值，为后续偏微分方程的数值求解提供了平滑且稳定的边界输入，保障了计算收敛性。

\subsubsection{时间推进与节点更新}

在每一个时间步内，获取下一时刻的环境值 \(T_\infty^{n+1}\) 和 \(C_\infty^{n+1}\) 后，对空间节点进行分类更新。其中，\(T_{N_r-2}^n\) 与 \(C_{N_r-2}^n\) 表示与表面相邻的内部节点在当前时刻的温度与水分浓度值。

\textbf{1. 内部节点更新（\(i = 1,2,\dots,N_r-2\)）}

温度与水分浓度的内部节点采用中心差分格式更新。由傅里叶导热定律和 Fick 扩散定律的柱坐标形式，可得温度更新公式：
\begin{equation}
T_i^{n+1} = T_i^n + \alpha \frac{\Delta t}{\Delta r^2} \left[ \left(1 + \frac{\Delta r}{2r_i}\right) T_{i+1}^n - 2T_i^n + \left(1 - \frac{\Delta r}{2r_i}\right) T_{i-1}^n \right]
\end{equation}
其中，热扩散系数 \(\alpha = k/(\rho c_p)\)，\(k\) 为药材热传导系数，\(\rho\) 为药材密度，\(c_p\) 为药材比热容。
同理，水分浓度更新公式为：
\begin{equation}
C_i^{n+1} = C_i^n + D_i^n \frac{\Delta t}{\Delta r^2} \left[ \left(1 + \frac{\Delta r}{2r_i}\right) C_{i+1}^n - 2C_i^n + \left(1 - \frac{\Delta r}{2r_i}\right) C_{i-1}^n \right]
\end{equation}
其中，\(D_i^n = D(C_i^n, T_i^n)\) 为当前时刻第 \(i\) 个节点的水分扩散系数，由题目附录中的经验公式计算得到。

\textbf{2. 中心节点更新（\(i = 0\)）}

在圆柱中心 \(r=0\) 处，由于对称性，热流与水分通量均为零。利用洛必达法则处理柱坐标方程的奇点，得到中心节点的更新公式：
\begin{equation}
T_0^{n+1} = T_0^n + 4\alpha \frac{\Delta t}{\Delta r^2} (T_1^n - T_0^n)
\end{equation}
\begin{equation}
C_0^{n+1} = C_0^n + 4D_0^n \frac{\Delta t}{\Delta r^2} (C_1^n - C_0^n)
\end{equation}
其中，\(D_0^n = D(C_0^n, T_0^n)\)。

\textbf{3. 表面节点更新（\(i = N_r-1\)）}

药材表面受到对流换热与对流传质的作用。记表面位置 \(r_{\text{surf}} = (N_r-1)\Delta r\)，引入曲率修正系数 \(A = 1 + \frac{\Delta r}{2r_{\text{surf}}}\)。结合第三类边界条件，表面温度更新公式为：
\begin{equation}
T_{N_r-1}^{n+1} = T_{N_r-1}^n + \frac{2\alpha\Delta t}{\Delta r^2}(T_{N_r-2}^n - T_{N_r-1}^n) - \frac{2\Delta t h}{\rho c_p \Delta r} A (T_{N_r-1}^n - T_\infty^{n+1})
\end{equation}
表面水分浓度更新公式为：
\begin{equation}
C_{N_r-1}^{n+1} = C_{N_r-1}^n + \frac{2D_{\text{surf}}^n \Delta t}{\Delta r^2}(C_{N_r-2}^n - C_{N_r-1}^n) - \frac{2\Delta t h_m}{\Delta r} A (C_{N_r-1}^n - C_\infty^{n+1})
\end{equation}
式中，\(h\) 和 \(h_m\) 分别为对流换热系数和对流传质系数；\(D_{\text{surf}}^n = D(C_{N_r-1}^n, T_{N_r-1}^n)\) 为表面节点处的水分扩散系数。此处，\(T_\infty^{n+1}\) 和 \(C_\infty^{n+1}\) 即为前文通过 PCHIP 插值获得的环境边界值。

\subsubsection{求解流程与结果输出}

上述显式差分的求解流程如图\ref{fig:algorithm_flowchart}所示。算法首先初始化温度与水分浓度矩阵，并计算热扩散系数；随后进入时间推进循环（\(n = 0\) 到 \(N_t-2\)），在每一时间步内，先利用 PCHIP 插值获取下一时刻的环境值 \(T_\infty^{n+1}\) 和 \(C_\infty^{n+1}\)，接着计算当前时刻各节点的水分扩散系数 \(D_i^n\)，并依次更新内部节点、中心节点及表面节点的温度与水分浓度；循环结束后，输出完整的温度矩阵 \(T_i^n\) 和水分矩阵 \(C_i^n\)，并按要求提取特定时刻与径向位置的切片数据。

值得注意的是，显式差分格式需满足 CFL 稳定性条件，即时间步长 \(\Delta t\) 必须足够小（一般要求 \(\alpha \frac{\Delta t}{\Delta r^2} \le \frac{1}{4}\)），以保证数值求解的收敛性。本文在后续计算中将严格选取满足该条件的时空步长。


\begin{figure}
    \centering
    \includegraphics[width=0.5\linewidth]{figures/algorithm_flow.pdf}
    \caption{一维非稳态热质耦合传热传质求解算法流程图}
    \label{fig:placeholder}
\end{figure}

\subsubsection{求解结果}






\subsection{数值算法的守恒性验证}

为验证上述显式有限差分算法求解结果的可靠性，本文分别从全局能量守恒和水分质量守恒两个维度，对数值解进行守恒性检验。

\subsubsection{能量守恒检验}

在连续域中，考虑半径 \(R\)、单位长度的圆柱状物料，其温度场满足柱坐标热传导方程：
\begin{equation}
\rho c_p \frac{\partial T}{\partial t} = \frac{1}{r} \frac{\partial}{\partial r} \left( k r \frac{\partial T}{\partial r} \right), \quad 0 < r < R
\end{equation}
表面 \(r=R\) 处为第三类对流换热边界：
\begin{equation}
-k \left. \frac{\partial T}{\partial r} \right|_{r=R} = h \left( T_\infty(t) - T(R,t) \right)
\end{equation}
将控制方程在 \([0,R]\) 上积分并代入边界条件，可得全局能量平衡方程。在 \([0,t]\) 上进一步积分，可推导得总能量变化等于边界热流累积：
\begin{equation}
\underbrace{Q(t) - Q(0)}_{\text{总能量变化}} = \underbrace{\int_0^t 2\pi R h \left( T_\infty(\tau) - T(R,\tau) \right) d\tau}_{\text{边界热流累积}}
\end{equation}
其中，总能量定义为 \(Q(t) = \int_0^R \rho c_p T(r,t) \cdot 2\pi r \, dr\)。

在离散域中，取节点 \(r_i = i \cdot \Delta r\)，对应控制体体积（单位长度）分别为 \(V_0 = \pi (\Delta r/2)^2\)，\(V_i = 2\pi r_i \Delta r\)（\(1 \le i \le N_r-2\)），\(V_{N_r-1} = \pi [R^2 - (R - \Delta r/2)^2]\)。
离散总能量计算式为：
\begin{equation}
Q^n = \sum_{i=0}^{N_r-1} \rho c_p T_i^n V_i
\end{equation}
边界热流累积为：
\begin{equation}
Q_{\text{flux}}^n = \sum_{m=0}^{n-1} h \cdot 2\pi R \left( T_\infty^m - T_{N_r-1}^m \right) \Delta t
\end{equation}
定义能量守恒相对误差为：
\begin{equation}
\varepsilon_Q^n = \frac{\left| (Q^n - Q^0) - Q_{\text{flux}}^n \right|}{\max \left( |Q^n - Q^0|, \epsilon \right)}
\end{equation}
其中，\(\epsilon\) 为防止分母为零的小量，本文取 \(\epsilon = 10^{-12}\)。

\subsubsection{水分质量守恒检验}

在连续域中，水分浓度场满足柱坐标变系数扩散方程：
\begin{equation}
\frac{\partial C}{\partial t} = \frac{1}{r} \frac{\partial}{\partial r} \left( r D(C) \frac{\partial C}{\partial r} \right), \quad 0 < r < R
\end{equation}
其中 \(D(C) = 7 \times 10^{-9} \exp\left(-\frac{0.89}{C}\right)\) 为随浓度变化的扩散系数。表面 \(r=R\) 处为第三类对流传质边界：
\begin{equation}
-D(C) \left. \frac{\partial C}{\partial r} \right|_{r=R} = h_m \left( C_\infty(t) - C(R,t) \right)
\end{equation}
同理，将控制方程积分并代入边界条件，可得总水分变化等于边界质流累积：
\begin{equation}
\underbrace{M(t) - M(0)}_{\text{总水分变化}} = \underbrace{\int_0^t 2\pi R h_m \left( C_\infty(\tau) - C(R,\tau) \right) d\tau}_{\text{边界质流累积}}
\end{equation}
其中，总水分定义为 \(M(t) = \int_0^R C(r,t) \cdot 2\pi r \, dr\)。

在离散域中，离散总水分计算式为 \(M^n = \sum_{i=0}^{N_r-1} C_i^n V_i\)，边界质流累积为 \(M_{\text{flux}}^n = \sum_{m=0}^{n-1} m^m \Delta t\)，其中表面质流 \(m^n = h_m \cdot 2\pi R \left( C_\infty^n - C_{N_r-1}^n \right)\)。同理可定义水分守恒相对误差 \(\varepsilon_M^n\)。

\subsubsection{变系数扩散的守恒性说明}

水分方程求解的关键难点在于 \(D(C)\) 是变量。若直接将 \(D\) 取节点值 \(D(C_i)\) 进行常规差分，则会漏掉 \(\partial D / \partial r \cdot \partial C / \partial r\) 项，导致相邻控制体界面通量不相等，从而破坏总水分守恒。

为克服这一问题，本文采用有限体积法思想，界面扩散系数取左右节点算术平均：
\begin{equation}
D_{i+1/2} = D\left( \frac{C_i + C_{i+1}}{2} \right)
\end{equation}
界面通量表示为：
\begin{equation}
F_{i+1/2} = D_{i+1/2} \frac{C_{i+1} - C_i}{\Delta r}
\end{equation}
此时，相邻控制体界面通量大小相等、方向相反，求和时自动抵消，从而在数学上保证了总水分严格守恒。

\subsubsection{守恒性检验结果}



采用时间步长 \(\Delta t = 1\,\text{s}\)，空间步长 \(\Delta r = 1\,\text{mm}\)，模拟总时长 \(t = 1800\,\text{s}\)。两项守恒检查结果如表 \ref{tab:conservation_check} 所示。

\begin{table}[H]
\centering
\caption{守恒性检验结果}
\label{tab:conservation_check}
\begin{tabular}{cccc}
\toprule
守恒量 & 最终相对误差 & 判断标准 & 结论 \\
\midrule
能量 & \(0.2349\%\) & \(< 1\% \sim 5\%\) & 优秀 \\
水分 & \(0.0004\%\) & \(< 1\% \sim 5\%\) & 优秀 \\
\bottomrule
\end{tabular}
\end{table}

结果表明，总能量变化曲线与边界热流累积曲线在整个模拟时段内几乎完全重合；总水分变化曲线与边界质流累积曲线同样几乎完全重合。两项守恒误差均远小于 5\% 的工程容许阈值，充分证明了本文所建立的显式有限差分算法具有良好的守恒性与数值稳定性，能够为后续中药材烘干过程的温度场与水分场求解提供可靠的数值保障。


%%%%%%%%%%%%%%%%%%%%%%%%%%%%%%%%%%%%%%%%%%%%%%%%%%%%%%%%%%%%% 
\newpage
\section{问题二的模型的建立和求解}

\subsection{模型建立：变物性湿热耦合控制方程}

问题一在建模时将密度 \(\rho\)、比热容 \(c_p\)、热传导系数 \(k\) 与水分扩散系数 \(D\) 视为常数。然而，在为期 \(2\sim3\) 天的整个烘干过程中，药材内部的水分浓度会发生显著变化，上述物性参数也随之改变，且烘房环境将由预热平衡阶段过渡到恒温干燥阶段。因此，问题二需在问题一的基础上，将物性参数改为随局部温度与水分浓度变化的函数，并考虑热质传递的双向耦合效应。

\subsubsection{两阶段烘房环境特征}

整个烘干过程可分为预热平衡与恒温干燥两个阶段，其环境特征分别为：
\begin{itemize}
    \item \textbf{预热平衡阶段}：烘房温度 \(T_\infty(t)\) 与水分浓度 \(C_\infty(t)\) 随时间明显变化；
    \item \textbf{恒温干燥阶段}：烘房温度 \(T_\infty(t)\) 接近恒定，水分浓度 \(C_\infty(t)\) 也进入平台。
\end{itemize}

\subsubsection{双向耦合关系}

在整个烘干过程中，药材内部温度与水分浓度之间存在双向耦合作用，其耦合路径可表示为：
\begin{align*}
    C &\rightarrow \rho, c_p, k \rightarrow T \\
    T, C &\rightarrow D \rightarrow C
\end{align*}
即水分浓度的变化通过改变密度、比热容与导热系数进而影响温度场，同时温度与水分浓度又共同决定水分扩散系数，进而反过来影响水分浓度场。

取圆柱中半径为 \(r\)、厚度为 \(dr\)、长度为 \(L\) 的环形控制体：
\begin{equation}
    dV = 2\pi r L \, dr
\end{equation}

\subsubsection{温度方程}

径向导热通量为：
\begin{equation}
    qr = -k(C)\frac{\partial T}{\partial r}
    \label{eq:qr}
\end{equation}

控制体内热量积累等于内表面流入减外表面流出：
\begin{equation}
    \rho(C) c_p(C) \, dV \frac{\partial T}{\partial t} = -\frac{\partial}{\partial r}(2\pi r L \cdot qr) \, dr
    \label{eq:energy_balance}
\end{equation}

将式 \eqref{eq:qr} 代入式 \eqref{eq:energy_balance}，约去 \(2\pi L \, dr\)，得到温度场控制方程：
\begin{equation}
    \rho(C) c_p(C) \frac{\partial T}{\partial t} = \frac{1}{r}\frac{\partial}{\partial r}\left[r k(C)\frac{\partial T}{\partial r}\right]
    \label{eq:temp_control}
\end{equation}

由附录3中的经验公式，展开有：
\begin{equation}
    \frac{1}{r}\frac{\partial}{\partial r}\left(r k \frac{\partial T}{\partial r}\right) = k\left(\frac{\partial^2 T}{\partial r^2} + \frac{1}{r}\frac{\partial T}{\partial r}\right) + k_c \frac{\partial C}{\partial r} \frac{\partial T}{\partial r}
    \label{eq:temp_expand}
\end{equation}
其中，\(k_c \frac{\partial C}{\partial r} \frac{\partial T}{\partial r}\) 表示水分分布对热传导的影响。

\subsubsection{水分方程}

等效水分通量为：
\begin{equation}
    qc = -D(C, T_k)\frac{\partial C}{\partial r}
    \label{eq:qc}
\end{equation}

对环形控制体进行水分守恒：
\begin{equation}
    dV \frac{\partial C}{\partial t} = -\frac{\partial}{\partial r}(2\pi r L qc) \, dr
    \label{eq:mass_balance}
\end{equation}

将式 \eqref{eq:qc} 代入式 \eqref{eq:mass_balance}，得到水分浓度场控制方程：
\begin{equation}
    \frac{\partial C}{\partial t} = \frac{1}{r}\frac{\partial}{\partial r}\left[r D(C,T) \frac{\partial C}{\partial r}\right]
    \label{eq:moisture_control}
\end{equation}

展开水分方程：
\begin{equation}
    \frac{\partial C}{\partial t} = D\left(\frac{\partial^2 C}{\partial r^2} + \frac{1}{r} \frac{\partial C}{\partial r}\right) + D_c \left(\frac{\partial C}{\partial r}\right)^2 + D_T \frac{\partial T}{\partial r} \frac{\partial C}{\partial r}
    \label{eq:moisture_expand}
\end{equation}

综上，问题二存在双向耦合：
\begin{align*}
    C &\rightarrow \rho, c_p, k \rightarrow T \\
    T, C &\rightarrow D \rightarrow C
\end{align*}

\subsection{初始条件与边界条件}

\subsubsection{初始状态}

烘干开始时，药材内部温度和水分浓度分布均匀，由题意可得
\begin{equation}
    T(r,0) = 28^\circ\text{C}, \qquad C(r,0) = 2.55 \, \text{kg/kg}, \qquad 0 \le r \le R
\end{equation}

\subsubsection{中心对称条件}

在圆柱中心 \(r=0\) 处，由轴对称性可知温度梯度与水分浓度梯度均为零，即
\begin{equation}
    \frac{\partial T}{\partial r}(0,t) = 0, \qquad \frac{\partial C}{\partial r}(0,t) = 0
\end{equation}

\subsubsection{表面 Robin 边界条件}

药材表面 \(r=R\) 处同时受到对流换热与对流传质作用，满足 Robin 边界条件：
\begin{align}
    -k(C_s) \frac{\partial T}{\partial r}(R,t) &= h \left[ T(R,t) - T_\infty(t) \right] \label{eq:bc_temp} \\
    -D(C(R,t), T_{R,s}) \frac{\partial C}{\partial r}(R,t) &= h_m \left[ C(R,t) - C_\infty(t) \right] \label{eq:bc_moisture}
\end{align}

其中，对流换热系数与对流传质系数分别为
\begin{equation}
    h = 25 \, \text{W/(m}^2\cdot\text{K)}, \qquad h_m = 8 \times 10^{-7} \, \text{m/s}
\end{equation}





\subsection{待删}
\textcolor{red}{\textbf{两阶段：}}
\begin{itemize}
    \item \textbf{预热平衡}：\(T_\infty(t), C_\infty(t)\) 明显随时间变化
    \item \textbf{恒温干燥}：\(T_\infty(t)\) 接近恒定，\(C_\infty(t)\) 也进入平台
\end{itemize}

\textbf{耦合关系：}
\begin{align*}
    C &\rightarrow \rho, c_p, k \rightarrow T \\
    T, C &\rightarrow D \rightarrow C
\end{align*}

取圆柱中半径为 \(r\)、厚度为 \(dr\)、长度为 \(L\) 的环形控制体：
\begin{equation}
    dV = 2\pi r L \, dr
\end{equation}

\subsection*{1. 温度方程}

径向导热通量为：
\begin{equation}
    qr = -k(C)\frac{\partial T}{\partial r}
\label{qr}
\end{equation}

控制体内热量积累等于内表面流入减外表面流出：
\begin{equation}
    \rho(C) c_p(C) \, dV \frac{\partial T}{\partial t} = -\frac{\partial}{\partial r}(2\pi r L \cdot qr) \, dr
\label{34}
\end{equation}

将\eqref{34}代入，约去 \(2\pi L \, dr\)，得到：
\begin{equation}
    \rho(C) c_p(C) \frac{\partial T}{\partial t} = \frac{1}{r}\frac{\partial}{\partial r}\left[r k(C)\frac{\partial T}{\partial r}\right]
\end{equation}

由附录3，展开有：
\begin{equation}
    \frac{1}{r}\frac{\partial}{\partial r}\left(r k \frac{\partial T}{\partial r}\right) = k\left(\frac{\partial^2 T}{\partial r^2} + \frac{1}{r}\frac{\partial T}{\partial r}\right) + k_c \frac{\partial C}{\partial r} \frac{\partial T}{\partial r} \quad \textcolor{red}{(存疑)}
\end{equation}
其中，\(k_c \frac{\partial C}{\partial r} \frac{\partial T}{\partial r}\) 表示水分分布对热传导的影响。

\subsection{2. 水分方程}

等效水分通量：
\begin{equation}
    qc = -D(C, T_k)\frac{\partial C}{\partial r}
\end{equation}

对环形控制体进行水分守恒：
\begin{equation}
    dV \frac{\partial C}{\partial t} = -\frac{\partial}{\partial r}(2\pi r L q_c) \, dr
\end{equation}

得到：
\begin{equation}
    \frac{\partial C}{\partial t} = \frac{1}{r}\frac{\partial}{\partial r}\left[r D(C,T) \frac{\partial C}{\partial r}\right]
\end{equation}

展开水分方程：
\begin{equation}
    \frac{\partial C}{\partial t} = D\left(\frac{\partial^2 C}{\partial r^2} + \frac{1}{r} \frac{\partial C}{\partial r}\right) + D_c \left(\frac{\partial C}{\partial r}\right)^2 + D_T \frac{\partial T}{\partial r} \frac{\partial C}{\partial r} \quad \textcolor{red}{(存疑)}
\end{equation}

因此 Q2 存在双向耦合：
\begin{align*}
    C &\rightarrow \rho, c_p, k \rightarrow T \\
    T, C &\rightarrow D \rightarrow C
\end{align*}

\subsection*{3. 初始与边界条件}

\textbf{初始状态为：}
\begin{equation}
    T(r,0) = 28^\circ\text{C}, \qquad C(r,0) = 2.55 \, \text{kg/kg}
\end{equation}

\textbf{中心满足轴对称条件：}
\begin{equation}
    \frac{\partial T}{\partial r}(0,t) = 0, \qquad \frac{\partial C}{\partial r}(0,t) = 0
\end{equation}

\textbf{表面 \(r=R\) 满足 Robin 边界：}
\begin{align*}
    -k(C_s) \frac{\partial T}{\partial r}(R,t) &= h \left[ T(R,t) - T_\textbf{air}(t) \right] \\
    -D(C(R,t), T_{R,s}) \frac{\partial C}{\partial r}(R,t) &= h_m \left[ C(R,t) - C_\textbf{air}(t) \right]
\end{align*}

其中，对流换热系数与对流传质系数分别为：
\begin{equation}
    h = 25 \, \text{W/(m}^2\cdot\text{K)}, \qquad h_m = 8 \times 10^{-7} \, \text{m/s}
\end{equation}

\subsection{求解结果}

%%%%%%%%%%%%%%%%%%%%%%%%%%%%%%%%%%%%%%%%%%%%%%%%%%%%%%%%%%%%% 

\section{问题三的模型的建立和求解}
\subsection{模型建立}

\subsection{模型求解}

\textbf{Step1：} 

\textbf{Step2：} 

\textbf{Step3：} 

\subsection{求解结果}

%%%%%%%%%%%%%%%%%%%%%%%%%%%%%%%%%%%%%%%%%%%%%%%%%%%%%%%%%%%%% 

\section{问题四的模型的建立和求解}

\subsection{模型建立：移动边界湿热耦合模型}

问题一至问题三均假设药材半径 \(R\) 为常数。然而，在实际烘干过程中，药材内部水分不断流失，会引起明显的体积收缩。此时，药材的计算区域不再是固定的 \(0 \le r \le R\)，而是随时间变化的移动边界区域 \(0 \le r \le R(t)\)。相应地，药材内部的材料点也会随着骨架发生径向位移。因此，问题四需在问题二变物性湿热耦合模型的基础上，进一步引入移动边界处理。

\subsubsection{收缩速度}

设药材半径 \(R(t)\) 随时间变化。假设药材沿径向均匀收缩，则位于半径 \(r\) 处的材料点收缩速度为：
\begin{equation}
v_r(r,t) = \frac{R'(t)}{R(t)} \cdot r
\end{equation}
式中，\(R'(t)\) 为半径随时间的变化率。该式表明，材料点的收缩速度与到中心的距离 \(r\) 成正比，越靠近表面的材料点移动越快，符合均匀径向收缩的物理图像。

\subsubsection{收缩材料中的守恒方程}

在发生收缩的材料中，物理量随某一固定材料点的变化率不能直接用空间偏导数描述，需引入随体导数（材料导数）：
\begin{equation}
\frac{D_m T}{Dt} = \frac{\partial T}{\partial t} + v_r \frac{\partial T}{\partial r}
\end{equation}
式中，\(\frac{\partial T}{\partial t}\) 为固定空间位置处的温度变化率；\(v_r \frac{\partial T}{\partial r}\) 为材料点移动引起的对流项。由能量与质量守恒定律，可得移动边界下的控制方程组：
\begin{equation}
\rho(C) c_p(C) \frac{D_m T}{Dt} = \frac{1}{r} \frac{\partial}{\partial r} \left[ r k(C) \frac{\partial T}{\partial r} \right]
\end{equation}
\begin{equation}
\frac{D_m C}{Dt} = \frac{1}{r} \frac{\partial}{\partial r} \left[ r D(C, T_k) \frac{\partial C}{\partial r} \right]
\end{equation}
方程组右端分别为径向导热项和径向扩散项。

\subsection{移动边界与无量纲坐标变换}

直接求解移动区域 \(0 \le r \le R(t)\) 上的偏微分方程，会导致空间网格随时间不断变化，计算复杂。为将移动边界问题转化为固定区域上的问题，本文引入无量纲坐标：
\begin{equation}
\xi = \frac{r}{R(t)}
\end{equation}
由此，物理域映射为固定计算域 \(0 \le \xi \le 1\)。其中，\(\xi=0\) 对应药材中心，\(\xi=1\) 对应药材当前表面。

令无量纲坐标下的温度与水分浓度分别为 \(\tilde{T}(\xi,t)\) 和 \(\tilde{C}(\xi,t)\)。由于 \(\xi\) 固定时表示同一材料点位置，因此沿材料点观察到的变化率等价于无量纲坐标下的时间偏导数：
\begin{equation}
\frac{D_m T}{Dt} = \frac{\partial \tilde{T}}{\partial t}, \qquad \frac{D_m C}{Dt} = \frac{\partial \tilde{C}}{\partial t}
\end{equation}

进一步，由链式法则推导空间偏导数的变换关系。由于空间求导时时间 \(t\) 固定，\(R(t)\) 对 \(r\) 而言为常数，故有：
\begin{equation}
\frac{\partial \xi}{\partial r} = \frac{1}{R(t)}
\end{equation}
若以温度为例，\(T(r,t) = \tilde{T}(\xi,t)\)，则：
\begin{equation}
\frac{\partial T}{\partial r} = \frac{\partial \tilde{T}}{\partial \xi} \frac{\partial \xi}{\partial r} = \frac{1}{R(t)} \frac{\partial \tilde{T}}{\partial \xi}
\end{equation}
因此，算子变换关系为：
\begin{equation}
\frac{\partial}{\partial r} = \frac{1}{R(t)} \frac{\partial}{\partial \xi}
\end{equation}
该变换表明，在同一时刻，关于实际距离 \(r\) 的变化率可通过关于归一化距离 \(\xi\) 的变化率表示。

\subsection{固定域上的控制方程}

将上述算子变换与随体导数关系代入移动边界下的守恒方程，可得固定域 \(0 \le \xi \le 1\) 上的控制方程组：
\begin{equation}
\rho(\tilde{C}) c_p(\tilde{C}) \frac{\partial \tilde{T}}{\partial t} = \frac{1}{R(t)^2 \xi} \frac{\partial}{\partial \xi} \left[ \xi k(\tilde{C}) \frac{\partial \tilde{T}}{\partial \xi} \right]
\end{equation}
\begin{equation}
\frac{\partial \tilde{C}}{\partial t} = \frac{1}{R(t)^2 \xi} \frac{\partial}{\partial \xi} \left[ \xi D(\tilde{C}, \tilde{T}_k) \frac{\partial \tilde{C}}{\partial \xi} \right]
\end{equation}
上述方程在形式上与固定半径情形一致，但引入了 \(1/R(t)^2\) 的尺度因子，体现了收缩对传热传质速度的影响。

\subsection{初始条件与边界条件}

\subsubsection{初始条件}
烘干开始时，药材内部温度和水分浓度分布均匀，由题意可得：
\begin{equation}
\tilde{T}(\xi,0) = 28^\circ\text{C}, \qquad \tilde{C}(\xi,0) = 2.55 \, \text{kg/kg}, \qquad 0 \le \xi \le 1
\end{equation}

\subsubsection{中心对称条件}
在药材中心 \(\xi=0\) 处，由轴对称性可知梯度为零：
\begin{equation}
\frac{\partial \tilde{T}}{\partial \xi}(0,t) = 0, \qquad \frac{\partial \tilde{C}}{\partial \xi}(0,t) = 0
\end{equation}

\subsubsection{表面热边界}
在当前表面 \(\xi=1\) 处，药材表面与烘房热风发生对流换热，满足 Robin 边界条件：
\begin{equation}
-\frac{k(\tilde{C}_s)}{R(t)} \frac{\partial \tilde{T}}{\partial \xi}(1,t) = h \left[ \tilde{T}_s - T_\infty(t) \right]
\end{equation}
其中，\(\tilde{T}_s = \tilde{T}(1,t)\) 表示药材表面温度，\(T_\infty(t)\) 为烘房空气温度。

\subsubsection{表面水分边界}
同理，在当前表面 \(\xi=1\) 处，药材表面与空气之间发生对流传质：
\begin{equation}
-\frac{D(\tilde{C}_s, \tilde{T}_{k,s})}{R(t)} \frac{\partial \tilde{C}}{\partial \xi}(1,t) = h_m \left[ \tilde{C}_s - C_\infty(t) \right]
\end{equation}
其中，\(\tilde{C}_s = \tilde{C}(1,t)\) 表示药材表面含水率，\(C_\infty(t)\) 为烘房空气水分浓度。

综上，问题四的模型在问题二的基础上，通过引入移动边界 \(R(t)\) 与无量纲坐标变换 \(\xi = r/R(t)\)，将原移动区域问题转化为固定域上的变系数偏微分方程组，从而可采用与问题二相同的显式有限差分法进行数值求解。求解时，\(R(t)\) 的取值由附件2中的药材半径随时间变化数据经插值给出。


\subsection{模型求解}

\textbf{Step1：} 

\textbf{Step2：} 

\textbf{Step3：} 

\subsection{求解结果}

%%%%%%%%%%%%%%%%%%%%%%%%%%%%%%%%%%%%%%%%%%%%%%%%%%%%%%%%%%%%%

\section{模型的分析与检验}

\subsection{结果检验}

\subsection{误差分析}

%%%%%%%%%%%%%%%%%%%%%%%%%%%%%%%%%%%%%%%%%%%%%%%%%%%%%%%%%%%%%

\section{模型的评价}

\subsection{模型的优点}
\begin{itemize}[itemindent=2em]
\item 优点1
\item 优点2
\item 优点3
\end{itemize}

\subsection{模型的推广与改进}

\subsection{模型的缺点}
\begin{itemize}[itemindent=2em]
\item 缺点1
\item 缺点2
\end{itemize}

%%%%%%%%%%%%%%%%%%%%%%%%%%%%%%%%%%%%%%%%%%%%%%%%%%%%%%%%%%%%%
%% 参考文献
\nocite{*}
\bibliographystyle{gbt7714-numerical}  % 引用格式
\bibliography{ref.bib}  % bib源

\newpage
%%%%%%%%%%%%%%%%%%%%%%%%%%%%%%%%%%%%%%%%%%%%%%%%%%%%%%%%%%%%%
%% 附录
\begin{appendices}
\section{文件列表}
\begin{table}[H]
\centering
\begin{tabularx}{\textwidth}{LL}
\toprule
文件名   & 功能描述 \\
\midrule
q1.m & 问题一程序代码 \\
q2.py & 问题二程序代码 \\
q3.c & 问题三程序代码 \\
q4.cpp & 问题四程序代码 \\
\bottomrule
\end{tabularx}
\label{tab:文件列表}
\end{table}

\section{代码}
\noindent q1.m
\lstinputlisting[language=matlab]{code/quesion1.m}
q2.py
\lstinputlisting[language=python]{code/q2.py}
q3.c
\lstinputlisting[language=c]{code/q3.c}
q4.cpp
\lstinputlisting[language=c++]{code/q4.cpp}
\end{appendices}
\end{document}


%%%%%双图模板%%%%%%
\begin{figure}
\centering
\subcaptionbox{炉温曲线示意图\label{fig:双图a}}
{\includegraphics[width=.4\textwidth]{炉温曲线示意图.png}}
\subcaptionbox{问题1炉温曲线\label{fig:双图b}}
{\includegraphics[width=.4\textwidth]{问题1炉温曲线.png}}
\caption{双图}\label{fig:双图}
\end{figure} 
%%%%%双图模板%%%%%%